% SGA 3, Expose VI_A, Loop 1 English reconstruction.
% Scope: paragraph 4.4 in full; paragraph 4.5 is excluded.
% Authority: Polo--Gille Exp6A-13oct24.pdf, component-local page 28
% (combined-reader page 331; printed page 320).
% OCR and the Reinhold English Markdown are comparison witnesses only.

\medskip
\phantomsection\label{par:sga3-VIA-4.4}%
\noindent\textbf{4.4.---}
还需构造当 \(G^\alpha\) 为 \(G\) 的任意连通分支时群胚
\(\mathscr S(G^\alpha_*)\) 的余核.
设 \(A'\) 是足够大的局部 \(A\)-代数,
它在 \(A\) 上有限且自由 (参见 \SGARefLink{thm:sga3-VIA-3.2}{3.2}).
则 \(G^\alpha\otimes_AA'\) 是 \(G\otimes_AA'\) 的有限个连通分支
\[
  C^1,\ldots,C^n
\]
的并, 且每个分支都有严格有理点.
因此, 对每个 \(i\), 存在 \(G\otimes_AA'\) 的右平移 \(r_i\),
将 \(G^0\otimes_AA'\) 映到 \(C^i\) 上. 这一平移诱导同构
\[
  \mathscr S(G^0_*)\otimes_AA'
    \xrightarrow{\ \sim\ }
  \mathscr S(C^i_*).
\]
所以, \(G_*\otimes_AA'\) 在 \(C^i\) 的饱和化上诱导的群胚有余核.

由于 \(\mathscr S(G^\alpha_*)\otimes_AA'\) 是某些群胚
\(\mathscr S(C^i_*)\) 的直和, 它也有余核.
这个余核是若干个
\[
  (F_0\otimes_AA')\backslash(G^0\otimes_AA').
\]
的副本的直和. 因此, 其中的每个有限子集都包含在某个仿射开子集中.
此外, 这个余核的形成与有限局部自由基扩张可交换
(参见编者注~\ref{note:sga3-VIA-37} 及下面的 \SGARefLink{par:sga3-VIA-4.6}{4.6}).
与 \SGARefLink{par:sga3-VIA-3.2.3}{3.2.3} 一样, 可知这个余核形如
\(Y\otimes_AA'\), 其中 \(Y\) 是 \(\mathscr S(G^\alpha_*)\) 的余核.
